% !TEX root = Projektdokumentation.tex

% Es werden nur die Abkürzungen aufgelistet, die mit \ac definiert und auch benutzt wurden.
%
% \acro{PBS}{Proxmox Backup Server}
% Ergibt in der Liste: PBS Proxmox Backup Server
% Im Text aber: \ac{PBS} -> Proxmox Backup Server (PBS)
%
% Hinweis: allgemein bekannte Abkürzungen wie z.B. bzw. u.a. müssen nicht ins Abkürzungsverzeichnis aufgenommen werden.
% Hinweis: allgemein bekannte IT-Begriffe wie Datenbank oder Programmiersprache müssen nicht erläutert werden,
%          aber ggfs. Fachbegriffe aus der Domäne des Prüflings.

% Die Option (in den eckigen Klammern) enthält das längste Label oder
% einen Platzhalter der die Breite der linken Spalte bestimmt.
\begin{acronym}[WWWWWW]
    \acro{DKFZ}{Deutsches Krebsforschungszentrum}
    \acro{ODCF}{Omics IT and Data Management Core Facility}
    \acro{PVE}{Proxmox Virtual Environment}
    \acro{PBS}{Proxmox Backup Server}
    \acro{VM}{Virtuelle Maschine}
    \acroplural{VM}[VMs]{Virtuelle Maschinen}
    \acro{CT}{Container}
    \acro{LXC}{Linux Container}
    \acroplural{LXC}[LXCs]{Linux Container}
    \acro{Ceph}{Ceph (verteiltes Storage-System)}
    \acro{RBD}{RADOS Block Device}
    \acro{LSF}{Load Sharing Facility}
    \acro{TSM}{IBM Tivoli Storage Manager (heute IBM Storage Protect)}
    \acro{VLAN}{Virtual Local Area Network}
    \acro{RAM}{Random Access Memory}
    \acro{RAID}{Redundant Array of Independent Disks}
    \acro{SFP}{Small Form-factor Pluggable}
    \acro{LACP}{Link Aggregation Control Protocol}
    \acro{GbE}{Gigabit Ethernet}
    \acro{API}{Application Programming Interface}
    \acro{SSH}{Secure Shell}
    \acro{LDAP}{Lightweight Directory Access Protocol}
    \acro{AD}{Active Directory}
    \acro{Prometheus}{Prometheus (Monitoring-System)}
    \acro{SOP}{Standard Operating Procedure}
    \acro{IHK}{Industrie- und Handelskammer}
\end{acronym}
